\chapter{Background and Related Work}
\label{ch:background}

This chapter establishes the concepts needed to assess the proposed system.
The literature is compared by sensing task, computational setting and
validation conditions. Results obtained on different datasets or devices are
not treated as a common numerical benchmark for this thesis.

\section{Visual fire and smoke detection}
\label{sec:background:vision}

RGB fire detection is a visual recognition task. A camera provides colour,
texture, boundary and motion cues, but it does not directly measure heat,
combustion gases or chemical products. Bright reflections, lamps, sunsets and
displayed fire imagery can resemble flame, whereas a distant flame may occupy
too few pixels to preserve useful structure. Smoke creates a different
appearance problem because it can be diffuse, semi-transparent and strongly
affected by background and illumination. A visual prediction must therefore
be described as evidence of a fire-like appearance rather than physical ground
truth.

UAV-oriented fire and smoke research has identified early small targets and
transparent smoke as difficult cases and has used multi-scale processing and
augmentation to address them \cite{Fan2025UAVFire}. That work supports the
need to examine target scale and environmental conditions. Its data, model and
test setting differ from those used here, so its reported accuracy is not a
baseline for the present system.

Domain shift is central to this application. Training images may contain large,
well-exposed outdoor fires, while a deployment camera observes small flames,
glare, clouds, reflective metal or indoor scenes. Closely related frames or
derived copies split between training and evaluation can also inflate a
performance estimate. Evidence from another imaging domain demonstrates how
subject-level leakage can produce optimistic deep-learning evaluation
\cite{Tampu2022Leakage}. The specific leakage mechanism differs here, but the
principle motivates source-family grouping, duplicate checks and frozen splits
in Chapter~4.

The final physical scope of this thesis is visible flame. Historical datasets
and models include smoke, but no controlled physical-smoke campaign was
conducted. Smoke-video playback can expose the detector to displayed imagery;
it cannot establish performance on airborne smoke. This distinction is needed
whenever fire and smoke are discussed together in the literature.

\section{Lightweight object detectors}
\label{sec:background:detectors}

Object detection assigns class scores and spatial boxes to candidate objects.
The original YOLO formulation cast this as one end-to-end regression from an
image to boxes and class probabilities \cite{Redmon2016YOLO}. The YOLOv8
software used in this project is a later engineering family. Ultralytics
documents its training, validation, inference and export interfaces, but no
formal YOLOv8 architecture paper exists \cite{UltralyticsYOLOv8}. The original
YOLO paper is therefore used for historical single-stage context and the
official software record identifies the implementation family; neither source
is used as performance evidence for the trained thesis checkpoints.

SSD predicts class scores and box offsets over default boxes at multiple
feature-map scales \cite{Liu2016SSD}. MobileNetV2 supplies a mobile-oriented
backbone based on inverted residual blocks and linear bottlenecks, and its
evaluation explicitly relates predictive quality to operations, parameter
count and latency \cite{Sandler2018MobileNetV2}. EfficientDet combines a
weighted bidirectional feature pyramid with compound scaling of resolution,
depth and width to address multiple resource budgets
\cite{Tan2020EfficientDet}. These design goals justify comparing SSD MobileNet
V2 FPNLite and EfficientDet-Lite0 with YOLOv8n and YOLOv8s. They do not predict
which trained implementation will perform best on the curated fire data or on
either edge device.

Input resolution and model capacity create coupled trade-offs. A larger input
can retain more pixels for a small object but increases the tensor and the work
performed by many layers. A larger model can represent more complex patterns
but increases memory, computation and often latency. Runtime libraries,
preprocessing, output decoding and accelerator support can reverse an ordering
suggested by parameter count alone. The candidate comparison in this thesis
therefore treats model family, size, resolution, exported graph and execution
backend as explicit configuration variables.

\section{Edge inference and model optimization}
\label{sec:background:edge}

Edge deployment separates at least three artifacts: a trained checkpoint, an
interchange or framework representation, and a runtime-specific optimized
artifact. Export success does not prove numerical fidelity, and a valid file
does not prove that the intended accelerator executes every operation. The
result must be checked against a frozen preprocessing and decoding contract and
then measured on the target device.

TensorRT constructs an optimized engine for a specified network, hardware and
precision configuration. NVIDIA documents reduced-precision execution and the
target-specific nature of serialized engines \cite{NvidiaTensorRT861}. In this
project, TensorRT FP16 is relevant to the Sentinel, while TensorFlow Lite is
used for the accepted Scout runtime. TensorFlow distinguishes float16 weight
conversion from full integer quantization; actual execution still depends on
the available delegates and kernels \cite{TensorFlowPostTrainingQuantization}.
An artifact labelled FP16 or quantized is therefore not, by itself, evidence
of accelerator use or end-to-end speed.

Latency also requires a declared boundary. \emph{Inference-call latency}
measures a model invocation after its input is available. \emph{Backend
pipeline latency} may add preprocessing and decoding. \emph{Camera-to-decision
latency} additionally includes capture, queues, temporal filtering,
communication, Scout processing and finalization. Frames per second computed
as the reciprocal of mean inference duration describe potential invocation
rate; they do not automatically equal camera or decision throughput. Chapters
6 and 7 preserve these distinctions.

The Jetson Nano and VOXL~2 also differ in architecture and software support.
NVIDIA specifies the Jetson Nano module with a quad-core Arm Cortex-A57 CPU, a
128-core Maxwell GPU and 4~GB LPDDR4 memory \cite{NvidiaJetsonNano}. ModalAI
specifies the VOXL~2 around the QRB5165 with eight CPU cores, 8~GB LPDDR5, an
Adreno~650 GPU and a neural-processing accelerator \cite{ModalAIVOXL2}.
Manufacturer specifications describe available resources; they do not predict
application latency, sustained temperature or power. Those quantities require
measurements of the exact artifact and concurrent services.

\section{Cooperative UAV perception and tiered verification}
\label{sec:background:cooperation}

Cooperative perception can exchange raw observations, intermediate features or
final decisions. These choices retain different amounts of information and
impose different bandwidth, synchronization and fusion requirements. Work on
adaptive feature fusion illustrates this early/intermediate/late distinction
in connected-vehicle LiDAR perception \cite{Qiao2023AdaptiveFusion}. Its
spatial scene-fusion results do not transfer to this camera-based two-drone
prototype, but the representation distinction is useful.

Collaborative UAV edge-intelligence proposals allocate computation and storage
across a swarm while considering node and link failures
\cite{Qu2024ElasticUAV}. Hierarchical aerial-computing research has also
examined cooperative resource allocation and task offloading
\cite{Kang2023HierarchicalAerial}. The present architecture addresses a
narrower problem. It does not partition a neural network or offload part of an
inference graph. Each vehicle executes a complete detector and exchanges a
compact decision-level message.

Decision-level exchange reduces communication volume and decouples the two
model runtimes, but it discards feature information that could support joint
fusion. It also requires the two cameras to observe a relevant scene at the
right time. In the evaluated stationary workflow, a Sentinel alert triggers a
full-frame Scout inference. The Sentinel box is retained as source-image
metadata; it is not projected into the Scout image or a world map. The Scout
therefore supplies a second visual decision but not geometric localization.

The tiered workflow assigns different intended priorities to its stages. The
Sentinel is intended to retain possible hazards for review, while the Scout is
intended to reject unsupported candidates before a final confirmation. Whether
this reduces false alerts while retaining true events depends on the data,
thresholds and error dependence between the two RGB models. The extra
verification and network path also adds latency and new failure modes. RQ3
tests this trade-off rather than assuming a benefit.

\section{Temporal filtering}
\label{sec:background:temporal}

A frame-wise detector can respond to a transient fire-like pattern that does
not persist. Video fire detection has therefore combined frame-level evidence
over time. Kong \emph{et al.} used temporal smoothing with a frame-level fire
classifier and reported a trade-off between sensitivity, delay and false
alarms \cite{Kong2016TemporalFire}. Their classifier, datasets and smoothing
method differ from this project, so their timing parameters and numerical
results are not transferred.

The mechanism studied here is a three-consecutive-frame rule. A qualifying
detection increments the current evidence count; loss of the qualifying event
resets it. An alert becomes eligible when the count reaches three. If the
per-frame events were independent and the probability of a false candidate in
one frame were \(p\), an idealized three-frame probability would be \(p^3\).
Real video frames are temporally correlated, however, so this expression is
not an estimator for the deployed system. Retained sequences and event-level
labels are required to measure suppression, retention and delay.

\section{ROS 2, DDS and correlated decisions}
\label{sec:background:ros2}

ROS~2 uses a middleware abstraction over DDS/RTPS implementations. It exposes
quality-of-service policies including history, depth, reliability, durability,
deadline and liveliness, and publisher/subscriber policies must be compatible
for communication \cite{OpenRoboticsROS2QoS}. ROS~2 can use different
middleware vendors, so implementation and compatibility are deployment
details rather than interchangeable labels \cite{OpenRoboticsRMW}.

A reliable DDS setting does not prove that the application produced one valid
answer within a deadline. The application must identify an event, reject
unrelated or late replies, define timeout behaviour and record a terminal
outcome. The protocol in this thesis consequently carries experiment, trial
and event identifiers and requires one accepted final decision for each
published alert.

Wireless ROS~2 studies with large payloads have identified fragmentation,
retransmission timing and buffer bursts as possible failure mechanisms
\cite{Lee2025ROS2Wireless}. Those results concern a different payload and
network setting. They motivate explicit message size, load and runtime
reporting but do not establish the behaviour of the compact SwarmNet event
messages used here. Likewise, MAVLink is designed for constrained robotic
links \cite{MAVLinkProtocolOverview}, but protocol design does not validate the
project's untested physical mission interface.

\section{Comparative synthesis and research gap}
\label{sec:background:gap}

Table~\ref{tab:related-work} summarizes how the principal sources inform this
thesis and why their validation settings cannot replace project-specific
experiments.

\begin{table}[H]
  \centering
  \caption{Comparison of related work and its role in this thesis. Reported
  findings remain conditional on each source's data, hardware and setting.}
  \label{tab:related-work}
  \scriptsize
  \begin{tabularx}{\textwidth}{@{}p{0.19\textwidth}p{0.18\textwidth}p{0.20\textwidth}X@{}}
    \toprule
    Source & Main topic & Validation setting & Use and limitation here \\
    \midrule
    Fan \emph{et al.} \cite{Fan2025UAVFire} & UAV fire/smoke detection & UAV
    remote-sensing imagery & Motivates small-target and smoke challenges;
    different data and model. \\
    Vazquez \emph{et al.} \cite{Vazquez2026EdgeFire} & Lightweight wildfire
    detection & Transfer-learned models with edge/resource evaluation &
    Motivates joint predictive/resource analysis; no numerical transfer. \\
    Liu \emph{et al.} \cite{Liu2016SSD}, Sandler \emph{et al.}
    \cite{Sandler2018MobileNetV2} & Single-stage and mobile detection & Generic
    object-detection benchmarks & Supports candidate design concepts; does not
    validate the exported SSD checkpoint. \\
    Tan \emph{et al.} \cite{Tan2020EfficientDet} & Scalable efficient detector &
    COCO object detection & Supports BiFPN/compound-scaling context; COCO results
    are not a fire-data baseline. \\
    Kong \emph{et al.} \cite{Kong2016TemporalFire} & Temporal fire evidence &
    Surveillance video & Motivates temporal processing; does not validate the
    three-frame rule. \\
    Qiao and Zulkernine \cite{Qiao2023AdaptiveFusion} & Cooperative perception &
    Connected-vehicle LiDAR fusion & Explains exchange levels; different sensor,
    geometry and task. \\
    Qu \emph{et al.} \cite{Qu2024ElasticUAV} & Collaborative UAV edge
    intelligence & Architectural resource/failure analysis & Motivates resource
    cooperation; unlike this decision-level two-node workflow. \\
    Kang \emph{et al.} \cite{Kang2023HierarchicalAerial} & Hierarchical aerial
    task offloading & Simulated learning-based resource allocation & Motivates
    hierarchical allocation; not a prototype baseline. \\
    ROS~2 documentation \cite{OpenRoboticsROS2QoS,OpenRoboticsRMW} & DDS QoS and
    middleware & Specified middleware semantics & Defines mechanisms; cannot
    establish SwarmNet latency or completion. \\
    \bottomrule
  \end{tabularx}
\end{table}

The literature establishes strong components: efficient single-stage
detectors, edge runtimes, temporal evidence, cooperative perception and typed
robotic communication. The research gap addressed here is their integration
into a traceable heterogeneous prototype in which model choice, artifact
identity, resource cost, temporal filtering, secondary verification and
mission logic are assessed under explicit evidence boundaries. The goal is not
to claim a universally novel swarm concept. It is to determine what this
particular division of responsibility supports, what it costs, and which
claims remain untested.
